\chapter{Fixed-point arithmetic}
\MARSChapterHeading{Fixed-point arithmetic}
\label{revision:fixed-point}

A pixel coordinate names a discrete cell, but an object can move by less than one
cell per update. If we discard that fraction after every addition, a velocity of
three quarters of a pixel can become no movement at all. The program needs a
position representation with finer steps than the display it eventually draws.

Fixed-point arithmetic supplies such a representation using integers. We choose
a scale, store scaled raw values, and keep that scale consistent through arithmetic.
The complete \texttt{fixed-point-ball.asm} program in \texttt{revision/labs/files}
uses signed words with eight fractional bits. It updates a point inside an
eight-by-eight frame, reflects it at the boundary, and converts the result to
one pixel only when rendering.

This chapter assumes signed words, shifts, function calls and row-major pixel
addressing. You should be able to explain a raw value's units, derive a reflected
position, choose a rounding rule, and justify an intermediate width. The independent
exercise interpolates between two signed raw values, including negative values
and full-word extremes that the motion model deliberately excludes.

\section{Choose a scale and keep the units}

With eight fractional bits, raw integer $r$ represents $r/256$. Raw 256 means one,
128 means one half, 192 means three quarters, and $-128$ means minus one half.
The stored bit pattern is still an ordinary two's-complement integer. The chosen
interpretation supplies the binary point; no register flag tells an addition
instruction which scale we intended.

For this motion model, position is measured in pixels and velocity in pixels per
update. Raw $x=320$ means 1.25 pixels. Raw $v_x=192$ means 0.75 pixels per update.
Adding those raw integers gives 512, representing two pixels after one update.
Both operands share a scale, so their sum keeps that scale. Adding literal one
would move by $1/256$ of a pixel; moving by one whole pixel requires adding 256.
An update is a simulation step, not a specified number of real seconds.

\begin{samepage}
  Every raw integer is an exact multiple of $1/256$ under this interpretation.
  Some desired values, such as one tenth, cannot be represented exactly at that
  scale. Choosing raw 26 approximates one tenth by 0.1015625; choosing raw 25 gives
  0.09765625. A rule for converting an external value is part of the interface.
  Using more fractional bits offers smaller steps but reduces the represented
  range available in a fixed-width word.
  \par
\end{samepage}

A signed 32-bit raw word therefore spans $-2^{31}/256$ through
$(2^{31}-1)/256$, in increments of $1/256$. This is the format's representable
range. It is much broader than the ball function's accepted range, and does not
promise that every arithmetic operation on two such words fits in another word.
Specify the range needed by each operation instead of relying on the format's name.

\section{Reflect an overshoot instead of clamping it}

The ball's raw walls are zero and $L=7\cdot256=1792$, the centers of display
columns or rows zero and seven. Horizontal coordinates increase rightward;
vertical coordinates increase downward. The axis function takes position
$0\leq p\leq1792$ and velocity $-1024\leq v\leq1024$. These are deliberate
preconditions, not checks for arbitrary external input.

First compute the trial $z=p+v$. If $0<z<L$, keep $z$ and the old velocity.
At or below the lower wall, return $-z$ and $-v$. At or above the upper wall,
return $2L-z$ and $-v$. Reflecting the excess keeps the portion of the step
traveled after contact. Clamping would discard that portion and change later states.

For example, raw 1728 is 6.75 pixels and raw 320 is 1.25 pixels per update.
Their trial is 2048, or eight pixels. It overshoots the upper wall by one pixel,
so the reflected position is six pixels: raw $3584-2048=1536$. The new velocity
is $-320$. Reversing velocity while leaving position at 2048 would still leave
the object outside the model's interval.

\begin{figure}[htbp]
  \centering
  \begin{tikzpicture}[x=12mm,y=9mm,alt={A one-dimensional interval runs from zero to seven pixels. The starting point is 6.75, raw 1728, moving right by 1.25, raw 320. Its trial is eight, raw 2048, one pixel beyond the upper wall. Reflecting that excess puts the result at six, raw 1536, with velocity minus 320.}]
    \draw[->] (-.2,0)--(8.5,0);
    \foreach \x in {0,...,8} {\draw (\x,-.1)--(\x,.1);\node[anchor=north,font=\small] at (\x,-.18) {\x};}
    \draw[thick] (7,-.15)--(7,2.5);
    \fill (6.75,0) circle (2pt);
    \draw[->,thick] (6.75,.55)--(8,.55);
    \draw[->,thick] (8,1.25)--(6,1.25);
    \node[anchor=south,font=\small] at (7.5,1.4) {reflect the excess};
    \node[anchor=west,font=\small] at (0,2.2) {$1728+320=2048$};
    \node[anchor=west,font=\small] at (0,1.5) {$3584-2048=1536$};
    \node[anchor=west,font=\small] at (0,.8) {new velocity $=-320$};
  \end{tikzpicture}
  \caption{Reflection preserves the traveled distance after crossing a wall. Coordinates on the axis are pixels; the equations use raw units of one two-hundred-fifty-sixth of a pixel.}
  \label{fig:revision-fixed-reflection}
\end{figure}


Exact contacts also reverse a nonzero incoming velocity. Position 1664 with
velocity 128 reaches 1792 exactly and returns velocity $-128$. Position 128
with velocity $-128$ reaches zero and returns velocity 128. Zero velocity stays
zero. An initial outward velocity at a wall is permitted; its first update
reflects inward under the same rule.

The complete axis routine has two results: position in \texttt{\$v0} and velocity
in \texttt{\$v1}. It uses only argument, result and temporary registers, leaving
saved registers and the stack pointer unchanged.

\noindent\begin{minipage}{\linewidth}
  \MARSCodeLabel{Reflect one bounded axis update}
  \begin{minted}{mars-idea}
advance_axis:
    addu $v0, $a0, $a1
    move $v1, $a1
    blezc $v0, reflect_low
    nop
    li $t0, 1792
    bgec $v0, $t0, reflect_high
    nop
    jrc $ra
reflect_low:
    subu $v0, $zero, $v0
    subu $v1, $zero, $v1
    jrc $ra
reflect_high:
    li $t0, 3584
    subu $v0, $t0, $v0
    subu $v1, $zero, $v1
    jrc $ra
\end{minted}
\end{minipage}

The inclusive lower and upper comparisons establish the contact convention.
The following \texttt{nop}s provide legal compact-branch fall-through slots;
they do not delay the simulated motion. A return ends each path with the pair
fully prepared.

The range proof is small enough to do explicitly. A trial lies between $-1024$
and 2816. Lower reflection produces zero through 1024; upper reflection produces
768 through 1792. The maximum speed is four pixels in a seven-pixel interval,
so one reflection suffices. All sums and negations fit in a signed word. This
proof does not cover a velocity of several wall-to-wall intervals or the most
negative word; either would require a broader contract and a different analysis.

\section{Round the display without changing the trajectory}

The model keeps raw coordinates between updates. Rendering converts a nonnegative
raw coordinate to a cell. Mode zero uses floor: $cell=\lfloor raw/256\rfloor$.
For these nonnegative inputs, a logical right shift by eight implements that
operation. Mode one uses nearest, with exact half ties upward:
$cell=\lfloor(raw+128)/256\rfloor$. Raw 127 maps to zero in either mode;
128 maps to zero in floor mode and one in nearest mode; 129 also maps to one
in nearest mode.

The renderer receives a full destination pointer in \texttt{\$a0}, raw $x$ and
$y$ in \texttt{\$a1}/\texttt{\$a2}, and mode zero or one in \texttt{\$a3}.
This fragment computes the returned pixel ID before the drawing loop:

\noindent\begin{minipage}{\linewidth}
  \MARSCodeLabel{Apply a display-only rounding bias}
  \begin{minted}{mars-idea}
render_ball:
    sll $t0, $a3, 7
    addu $t1, $a1, $t0
    addu $t2, $a2, $t0
    srl $t1, $t1, 8
    srl $t2, $t2, 8
    sll $t2, $t2, 3
    addu $v0, $t2, $t1
\end{minted}
\end{minipage}

Shifting mode by seven produces either zero or 128. The biased coordinates reach
at most 1920, so the additions remain safe. Each shifted cell remains zero through
seven. The row shift by three and final addition form $8\cdot row+column$.
The complete function writes all 64 destination words: one white pixel and 63
black pixels. It uses a full-width cursor increment for addresses, preserves saved
registers and \texttt{\$sp}, and returns the ID in \texttt{\$v0}.

The destination argument matters even though main always passes \texttt{frame}.
Independent callers can pass \texttt{spare\_frame} or another valid allocation.
The renderer must not quietly replace that pointer with the display's global address.
Nor may it replace the model coordinates with their rounded cells. The two display
modes should change only pixels, not subsequent motion.

\begin{table}[htbp]
\centering
\small
\begin{tabular}{rrrrrrr}
\toprule
Tick & $x$ raw & $y$ raw & $v_x$ raw & $v_y$ raw & Floor ID & Nearest ID\\
\midrule
0 & 320 & 448 & 192 & 320 & 9 & 17\\
1 & 512 & 768 & 192 & 320 & 26 & 26\\
2 & 704 & 1088 & 192 & 320 & 34 & 35\\
3 & 896 & 1408 & 192 & 320 & 43 & 52\\
4 & 1088 & 1728 & 192 & 320 & 52 & 60\\
5 & 1280 & 1536 & 192 & -320 & 53 & 53\\
6 & 1472 & 1216 & 192 & -320 & 37 & 46\\
7 & 1664 & 896 & 192 & -320 & 30 & 39\\
8 & 1728 & 576 & -192 & -320 & 22 & 23\\
\bottomrule
\end{tabular}
\caption{The default trajectory retains the same raw state in both display modes. Each ID is eight times the rounded row plus the rounded column. Tick zero is the initialized state, before any update.}
\label{tab:revision-fixed-trace}
\end{table}


At tick seven, $x=1664$ is 6.5 pixels, still moving right. Nearest mode displays
column seven already. At tick eight, reflection has occurred and the raw position
is 1728 with velocity $-192$; nearest mode still displays column seven. A pixel
alone cannot tell you whether contact has occurred. Use the raw position and
velocity to explain that state. Quantization is a many-to-one conversion, so
different model states can have the same displayed coordinate.

\section{Publish the initial state and each complete update}

Main loads the four seeds into saved registers, retains the requested count and
rounding mode, and starts tick at zero. The default seeds are raw
$(x,y,v_x,v_y)=(320,448,192,320)$. It writes those raw words and tick to the public
state, calls the renderer, stores its returned pixel ID, then reaches
\texttt{frame\_ready}. The single \texttt{nop} before \texttt{frame\_continue}
is a countable observation marker, not a timer.

If no updates remain, main exits. Otherwise it calls \texttt{advance\_axis} for
$x$, retains both results, then calls it for $y$ with its own position and velocity.
Only after both results are retained does it increment tick and publish again.
The saved registers protect model state from allowed temporary-register changes
inside either helper. Main eventually publishes ticks zero through $N$, so a
request of eight produces nine observations. A request of zero still initializes
and displays its seed.

A debugger paused between the two axis calls can see an intermediate state in
registers: new horizontal values and old vertical values. That is not a complete
published tick. At \texttt{frame\_ready}, raw words, tick, pixel ID and frame must
agree. The caller owns this publication order; the renderer owns only its returned
ID and destination pixels. Input seeds, options, guards and the unused spare frame
remain unchanged, including across stores that might otherwise be restored later.

\section{Choose intermediate width before reducing precision}

The bounded motion proof justified word arithmetic. The separate
\texttt{add\_fixed\_probe} accepts any two correctly sign-extended signed words
and exposes a case outside that proof. It returns the wrapped word sum in
\texttt{\$v0} and exact doubleword sum in \texttt{\$v1}. Its body uses
\texttt{addu} and \texttt{daddu} on the same operands.

For raw inputs hexadecimal \texttt{7FFFFF00} and \texttt{00000200}, the exact
raw sum is 2,147,483,904, hexadecimal \texttt{0000000080000100}. It does not fit
in a signed word. \texttt{addu} keeps low word \texttt{80000100} and sign extends
it to hexadecimal \texttt{FFFFFFFF80000100}, interpreted as $-2{,}147{,}483{,}392$.
\texttt{daddu} retains the exact positive result because the sum fits in 64 bits.
Both still describe raw units at the original scale; changing width did not move
the binary point.

The unsigned suffix in \texttt{addu} specifies nontrapping word arithmetic; it
does not make the result a nonnegative doubleword. For this target, word results
are sign extended. \texttt{add} instead raises a signed-overflow exception for
this pair. Set \texttt{DEMO} to one in a copy and stop at
\texttt{overflow\_ready} to compare the results before exit overwrites
\texttt{\$v0}. A temporary \texttt{add} substitution investigates the exception;
restore the working program afterward.

Multiplication changes scale as well as range. If $a$ and $b$ represent values
$a/256$ and $b/256$, their integer product represents a quantity scaled by
$256^2$. Returning to eight fractional bits requires dividing that product by
256 with an explicit rounding rule. Raw $192\cdot320=61440$ reduces to raw 240,
representing $0.9375$. Narrowing the product before rescaling can lose bits that
the final result needs.

\begin{samepage}
  The Release 6 \texttt{dmul} instruction multiplies doubleword operands and keeps
  the low 64 bits in an ordinary result register. It is exact only when the needed
  product fits in that width. Two signed-word operands have a signed product that
  fits in 64 bits; unrestricted doubleword operands do not have that guarantee.
  A signed \texttt{dsra} shift retains the sign and rounds division by a power of
  two downward. A logical shift would turn a negative intermediate into a large
  positive value. These width and rounding decisions become part of the next
  exercise's contract.
  \par
\end{samepage}

\section{Interpolate signed raw values independently}

Implement \texttt{interpolate\_raw} in \texttt{fixed-interpolate.asm}. It takes
two correctly sign-extended signed raw words, $a$ and $b$, and an integer weight
$w$ from zero through 256. Return the nearest raw integer to
$((256-w)a+wb)/256$, with exact half ties toward positive infinity. All inputs
share the same eight-fractional-bit format. Weight zero selects $a$, weight 256
selects $b$, and weight 128 selects their midpoint.

For raw $a=-256$, $b=512$, $w=64$, the weighted numerator is $-16384$ and the
returned raw value is $-64$, representing minus one quarter. This interpolates
one quarter of the way from minus one to two. Do not quantize either endpoint
to a display cell before computing it.

The coefficients are nonnegative and sum to 256. The exact numerator is bounded
by $-2^{39}$ and $2^{39}-256$, so a signed doubleword has ample range even though
a word does not. The returned rounded raw value stays between the endpoints and
fits in a signed word. If you choose the algebraic form $a+w(b-a)/256$, widen
before subtracting: two valid endpoints can have a difference that exceeds a
signed word. Establish bounds for every intermediate in your chosen form.

Use concrete negative half cases to establish rounding: the nearest integer to
$-0.5$ with this tie rule is zero, while the nearest integer to $-1.5$ is $-1$.
Adding a half-unit bias to a wide numerator and then taking floor can implement
this convention, but a division that truncates toward zero does not implement
floor for negative values. Preserve saved registers, \texttt{\$sp}, caller-owned
stack bytes, inputs and guards. Choose your temporary registers and algorithm;
the written lab supplies public checks and boundary investigations.
